% TOP_PROBLEM: {"id": 30006937, "title": "LeBrun-Salamon Conjecture: positive quaternion-Kahler manifolds are symmetric spaces", "category_id": 6, "category": {"id": 6, "name": "geometry", "display_name": "Geometry", "description": "Euclidean and non-Euclidean geometry, geometric structures.", "slug": "geometry", "order_index": 6, "created_at": "2026-07-31T15:26:25.671Z"}, "status": "open"}
% FIRST_POSTED: 2026-09-27
% ATTEMPT_STATUS: unresolved
% !TeX program = lualatex
\documentclass[11pt]{article}
\usepackage[margin=25mm]{geometry}
\usepackage{fontspec}
\setmainfont{Latin Modern Roman}
\usepackage{amsmath,amssymb,mathtools}
\usepackage{unicode-math}
\setmathfont{Latin Modern Math}
\usepackage{xurl,hyperref}
\hypersetup{colorlinks=true,urlcolor=blue}
\setcounter{secnumdepth}{0}
\setlength{\parindent}{0pt}
\setlength{\parskip}{5pt}
\setlength{\emergencystretch}{3em}
\begin{document}
\section*{\texorpdfstring{LeBrun-Salamon Conjecture: positive quaternion-Kahler manifolds are symmetric spaces}{LeBrun-Salamon Conjecture: positive quaternion-Kahler manifolds are symmetric spaces}}\label{30006937}
\subsection{Definitions and mathematical statement}
In real dimension $4n$, $n\ge2$, a quaternion-K\"ahler manifold is a Riemannian manifold whose holonomy lies in $\operatorname{Sp}(n)\operatorname{Sp}(1)$. Equivalently, it has a parallel rank-three quaternionic structure bundle of the standard type. Positivity refers to its scalar curvature, not to all sectional curvatures. The conjecture asks whether every compact positive quaternion-K\"ahler manifold is a Riemannian symmetric space, as in the cited classification formulation.

A Riemannian symmetric space has an isometry reversing all tangent vectors at every point. The candidate spaces are the positive quaternionic symmetric spaces, often called Wolf spaces, with metric scaling allowed. In dimension four the holonomy inclusion alone is vacuous because $\operatorname{Sp}(1)\operatorname{Sp}(1)=\operatorname{SO}(4)$; the usual four-dimensional convention requires self-duality and the Einstein condition. Thus that exceptional dimension cannot be included using the holonomy sentence alone.
\par\smallskip\textit{Formulation and concept references:} \href{https://link.springer.com/article/10.1007/s12220-024-01624-7}{[S1]}.

\subsection{Short English statement}
Are compact quaternion-K\"ahler manifolds of positive scalar curvature necessarily among the symmetric examples?
\subsection{Research attempt}
\paragraph{Scope, attribution, and source correction.}
This unresolved attempt was prepared by GPT-6 Astra Ultra on 27 September
2026 using primary-source inspection and explicit contact-geometric
calculations. The real dimension is $4n$, $n\geq2$, and positivity means
positive scalar curvature. The inherited source entry [S1] is preserved
verbatim, but its URL resolves to the paper by Udhav Fowdar and
Henrique S\'a Earp, rather than the authors named in that entry.
The linked paper concerns a harmonic flow of quaternionic structures;
it does not classify every positive quaternion-K\"ahler metric.
For the geometric reduction we use LeBrun's primary paper [E1].
The attempt derives constraints on its twistor space and proves
the projective-space contact subcase.

\paragraph{The established twistor input.}
A positive quaternion-K\"ahler manifold has a twistor space
$\pi:Z\to M$ with fiber $\mathbb P^1$, complex dimension $2n+1$,
a holomorphic contact structure and a positive K\"ahler--Einstein
metric. These are established geometric results used here.
Conversely the contact K\"ahler--Einstein construction in [E1]
recovers a positive quaternion-K\"ahler manifold. Thus a
classification of the relevant contact K\"ahler--Einstein
spaces, with their geometric reconstruction, would give a
classification route for the original metric problem.

This is a reduction, not a claim that every Fano manifold
is homogeneous or that every complex contact manifold
automatically carries such a metric. The K\"ahler--Einstein
and contact hypotheses must travel together. The original
target also concerns the Riemannian metric, so identifying
a smooth manifold alone would be insufficient.

\paragraph{A determinant identity from the contact form.}
Let $D\subset TZ$ be the rank-$2n$ contact distribution and
$L=TZ/D$ its quotient line bundle. Locally the quotient
map is an $L$-valued one-form $\theta$. Although $d\theta$
depends on a trivialization of $L$, its restriction to
$D$ does not: under $\theta\mapsto f\theta$ the term
$df\wedge\theta$ vanishes on $D$. Contact nondegeneracy
says that $(d\theta)|_D$ is a nondegenerate alternating
form with values in $L$.

Taking its $n$-th exterior power gives a nowhere-zero
section of $(\det D)^*\otimes L^n$, hence
$\det D\cong L^n$. The exact sequence
$0\to D\to TZ\to L\to0$ now gives
\begin{equation}\label{a442:canonical}
                     K_Z^{-1}\cong L^{n+1}.
\end{equation}
This is an isomorphism of line bundles, not merely
an equality of real Chern classes with possible
torsion suppressed. Equivalently
$\theta\wedge(d\theta)^n$ trivializes
$K_Z\otimes L^{n+1}$.

For a twistor space, $K_Z^{-1}$ is ample, so $L$ is
ample too: a positive tensor power of a line bundle
is ample precisely when the line bundle is ample.
For every curve $C$,
$(-K_Z)\cdot C=(n+1)L\cdot C$.
These divisibility and positivity restrictions are
necessary, but they do not classify $Z$. Many
integer intersection data can satisfy a divisibility
condition without determining a variety.

\paragraph{A complete algebraic subcase: projective twistor space.}
Assume additionally that $Z\cong\mathbb P^{2n+1}$.
Its Picard group is generated by $\mathcal O(1)$ and
$K_Z^{-1}=\mathcal O(2n+2)$. Equation
\eqref{a442:canonical} forces $L=\mathcal O(2)$.
The Euler sequence, twisted by $\mathcal O(1)$,
identifies
\[
 H^0(\mathbb P(V),\Omega^1(2))
 =\ker\bigl(V^*\otimes V^*\longrightarrow
                                  \operatorname{Sym}^2V^*\bigr)
 =\Lambda^2V^*,
 \qquad\dim V=2n+2.
\]
Thus a contact quotient form is represented by an
alternating bilinear form $A$ on $V$.

If $A$ has a nonzero radical vector $z$, the quotient
form vanishes at $[z]$ and cannot define a contact
distribution there. If $A$ is nondegenerate, its
induced form on $z^{\perp_A}/\mathbb Cz$ is
nondegenerate; this is precisely the contact
nondegeneracy on the rank-$2n$ distribution.
Therefore the contact forms are exactly the
nondegenerate alternating forms.

All such forms on a fixed even-dimensional complex
vector space are equivalent under $\operatorname{GL}(V)$,
by the elementary symplectic-basis construction.
Consequently all contact structures of this form
on projective space are equivalent under projective
automorphisms. Combining this explicit classification
with the established twistor reconstruction and
uniqueness up to metric scale gives the standard
quaternionic projective space $\mathbb H P^n$.
The additional assumption $Z\cong\mathbb P^{2n+1}$
is the entire restriction; it has not been deduced
for general $M$, and the other Wolf spaces would
not satisfy it.

\paragraph{A topological check from the twistor fibration.}
The contact line restricts to $\mathcal O(2)$ on
each twistor line. Thus
$h=\tfrac12c_1(L)\in H^2(Z,\mathbb Q)$ restricts
to the generator of $H^2(\mathbb P^1,\mathbb Q)$.
The classes $1,h$ restrict to a basis of the
fiber cohomology. Leray--Hirsch therefore gives,
as graded rational vector spaces,
\[
 H^k(Z,\mathbb Q)\cong
 H^k(M,\mathbb Q)\oplus H^{k-2}(M,\mathbb Q).
\]
In particular $b_2(Z)=b_2(M)+1$. This argument
uses an actual global class restricting to the
fiber generator; it does not simply assume that
every sphere bundle is a cohomological product.

For the projective case, $b_2(Z)=1$, so
$b_2(M)=0$, consistent with $\mathbb H P^n$.
The converse inference would be invalid:
$b_2(M)=0$ gives $b_2(Z)=1$, not an
identification of $Z$ with projective space.
Picard number one is much weaker than a
complete holomorphic classification. Nor
does the vector-space decomposition assert
a product ring structure or a topologically
trivial fibration.

\paragraph{Why immediate flow and curvature arguments stop.}
Let $\Omega$ be the parallel fundamental quaternionic
four-form of an actual quaternion-K\"ahler metric.
A torsion energy of the form
$\tfrac12\int_M|\nabla\Omega|^2$ is zero.
Its negative gradient therefore vanishes at
this structure: a differentiable nonnegative
functional has zero derivative at a zero minimum.
Running such a flow starting from the unknown
torsion-free structure does not move it toward
a symmetric example. One still has to classify
the zero-energy structures. The linked harmonic
flow paper also distinguishes harmonic structures
from torsion-free ones; merely reaching a
critical point would not resolve that issue.

Positive quaternion-K\"ahler metrics are Einstein,
so one may normalize
$\operatorname{Ric}=(4n-1)g$. Myers' theorem then
bounds the diameter by $\pi$. This is useful
geometric control, but it is not a bound on the
full curvature tensor or a proof that
$\nabla\operatorname{Rm}=0$. The latter is the
local symmetry condition that a classification
argument would need to establish or replace.
Even among the symmetric candidate spaces,
higher rank permits zero-curvature planes;
positive scalar curvature must not be
strengthened to positive sectional curvature.

Thus an Einstein normalization, a diameter
bound and the contact Chern relation are
compatible necessary restrictions, but none
eliminates all hypothetical nonsymmetric
examples. Compactness or finite lists of
topological data alone would still require
a rigidity argument for the metric.

\paragraph{Verification, first gap, and outcome.}
The contact exponent $n+1$ was checked from the
rank-$2n$ distribution, and the projective
calculation gives $\mathcal O(2)$ rather than
$\mathcal O(1)$. The rational fiber generator
retains its factor of one half. Metric scale
is allowed throughout, and the exceptional
four-dimensional holonomy convention is not
used to enlarge the target.

The first remaining classification input is to
show that every twistor contact K\"ahler--Einstein
space in this scope is one of the homogeneous
models, or to construct a valid nonhomogeneous
one with its required reconstruction.
\textbf{UNRESOLVED.} The determinant identity,
projective contact classification and topology
check are proved. They do not classify all
positive quaternion-K\"ahler manifolds, and
no novelty is claimed.

\subsection{Sources}
\begin{itemize}
\item[S1]Fadel, Loubeau, and Moreno, Harmonic Flow of Quaternion-Kahler Structures, 2024, Introduction.
\url{https://link.springer.com/article/10.1007/s12220-024-01624-7}.
\textit{Formulation source.}
\item[E1]Claude LeBrun, \emph{Fano Manifolds, Contact Structures, and Quaternionic Geometry}, International Journal of Mathematics 6 (1995), 419--437.
\url{https://www.math.stonybrook.edu/~claude/fcq.pdf}.
\textit{Primary author text for the twistor/contact K\"ahler--Einstein correspondence; consulted 27 September 2026. The inherited [S1] author mismatch is identified in the attempt rather than silently changing its source entry.}
\end{itemize}
\end{document}
